% ECONOMICS_PROBLEM: {"id": "C78-217", "title": "Lattice structure under one preference change on each side", "jel_code": "C78", "source_id": "OP-0200"}
% FIRST_POSTED: 2026-10-10
% ATTEMPT_STATUS: solved
% ENTRY_KIND: research_attempt
\documentclass[11pt]{article}
\usepackage[margin=1in]{geometry}
\usepackage{amssymb,amsmath,amsthm,enumitem}
\usepackage[hidelinks]{hyperref}

\newtheorem{theorem}{Theorem}
\newtheorem{lemma}[theorem]{Lemma}
\newtheorem{corollary}[theorem]{Corollary}
\newtheorem{proposition}[theorem]{Proposition}
\theoremstyle{definition}
\newtheorem{definition}[theorem]{Definition}
\newtheorem{example}[theorem]{Example}
\theoremstyle{remark}
\newtheorem{remark}[theorem]{Remark}

\newcommand{\D}{\mathcal D}
\newcommand{\Sol}{\operatorname{Sol}}
\newcommand{\cur}{\operatorname{cur}}
\newcommand{\nxt}{\operatorname{next}}

\title{One preference change on each side: which robust sublattices occur}
\author{Claude Opus 5.5 (Anthropic), submitted by Ji Ho Bae}
\date{October 10, 2026}

\begin{document}

\maketitle
\begin{abstract}
Let two stable-marriage profiles $A,B$ differ in at most one worker's list and at most one firm's list. Put $L=\mathcal M_A$ and $R=\mathcal M_A\cap\mathcal M_B$. We determine exactly which pairs $(L,R)$ occur, up to lattice isomorphism. Writing $L$ as the down-set lattice of its join-irreducibles $P$, the occurring sets $R$ are exactly those cut out by two chains of $P$: an arbitrary set of admissible positions on each chain, and implications whose tail lies on the first chain or whose head lies on the second. In particular, not every sublattice occurs. The complement $L\setminus R$ can be closed under meets, under joins, under both, or under neither.
\end{abstract}

\section{The answer}

Throughout, $W$ and $F$ are equal-sized finite sets of workers and firms, and profiles are strict and complete. A matching $M$ is stable for a profile if no worker--firm pair strictly prefers each other to their partners. The stable matchings $\mathcal M_A$ form a distributive lattice under unanimous weak worker preference, the matching better for workers being smaller.

For a finite poset $P$ let $\D(P)$ be the lattice of its down-sets under inclusion. Birkhoff's theorem gives $L\cong\D(J(L))$, where $J(L)$ is the poset of join-irreducible elements.

\begin{definition}\label{def:system}
A \emph{one-change system} on $P$ consists of chains $C_1=(\alpha_1<\dots<\alpha_p)$ and $C_2=(\beta_1<\dots<\beta_q)$ of $P$ (either may be empty), sets $A_1\subseteq\{0,\dots,p\}$ and $A_2\subseteq\{0,\dots,q\}$, and a finite set $\Phi$ of implications $s\to t$ of two kinds:
\[
 \text{($w$-kind)}\quad s\in C_1\cup\{\top\},\ t\in P\cup\{\bot\};\qquad
 \text{($f$-kind)}\quad s\in P\cup\{\top\},\ t\in C_2\cup\{\bot\}.
\]
An implication can be of both kinds. A down-set $X$ \emph{satisfies} $s\to t$ if $s\in X$ implies $t\in X$, where $\top\in X$ and $\bot\notin X$ by convention. The solution set is
\[
 \Sol=\{X\in\D(P):\ |X\cap C_1|\in A_1,\ |X\cap C_2|\in A_2,\ X\text{ satisfies every member of }\Phi\}.
\]
\end{definition}

The maps $X\mapsto|X\cap C_i|$ are lattice homomorphisms onto chains, and each implication defines a sublattice, so every solution set is a sublattice of $\D(P)$ or empty.

\begin{theorem}\label{thm:main}
Let $L$ be a finite distributive lattice, $P=J(L)$, and $R\subseteq L$. There are profiles $A,B$, differing in at most one worker's list and at most one firm's list, and a lattice isomorphism $\mathcal M_A\to L$ carrying $\mathcal M_A\cap\mathcal M_B$ onto $R$, if and only if $R$ corresponds, under $L\cong\D(P)$, to the solution set of a one-change system on $P$.
\end{theorem}

\begin{example}\label{ex:no}
Let $P=\{x_1,x_2,x_3,x_4\}$ be an antichain, so $L\cong\{0,1\}^4$, and let $R=\{X:\ x_1\in X\Leftrightarrow x_2\in X,\ x_3\in X\Leftrightarrow x_4\in X\}$. This is a sublattice, but it does not occur. Suppose a system has $\Sol=R$. Every chain of $P$ has at most one element, and $R$ contains $\varnothing$ and $P$, so $A_1$ and $A_2$ contain every position and exclude nothing. The only nontrivial implications valid on $R$ are $x_1\to x_2$, $x_2\to x_1$, $x_3\to x_4$ and $x_4\to x_3$, and each is needed: the down-set $\{x_1\}$ violates only the first, and so on. An admissible implication has its tail in $C_1$ or its head in $C_2$. One element of $C_1$ is the tail of only one of the four, and one element of $C_2$ is the head of only one. So at most two are admissible, and $\Sol\ne R$.
\end{example}

\begin{corollary}\label{cor:complement}
Put $D=L\setminus R$. Among the pairs that occur, $D$ can be meet-closed or not and join-closed or not, in all four combinations. So these properties are not forced by realizability and must be checked on $(L,R)$ directly.
\end{corollary}

\begin{proof}
Take $P=\{x,y\}$ an antichain and $C_1=C_2=\{x\}$, with all positions allowed. The implications $x\to y$ (tail in $C_1$) and $y\to x$ (head in $C_2$) give $R=\{\varnothing,\{x,y\}\}$, whose complement is closed under neither operation. The implications $\top\to x$ and $\top\to y$ give $R=\{\{x,y\}\}$, whose complement is meet-closed but not join-closed. The implications $x\to\bot$ and $y\to\bot$ give $R=\{\varnothing\}$, whose complement is join-closed but not meet-closed. Finally $P=\{x\}$ with $C_1=C_2=\{x\}$ and $\top\to x$ gives a complement $\{\varnothing\}$ closed under both.
\end{proof}

\section{Rotations}

We use the rotation theory of Gusfield and Irving \cite[Ch.~2--3]{GI}. Fix a profile $A$, and let $M_0$ be its worker-optimal stable matching, which worker-proposing deferred acceptance returns \cite{GS}. For a stable $M$ and a worker $x$, the successor $s_M(x)$ is the first firm after $M(x)$ on $x$'s list that prefers $x$ to its partner, if one exists; put $\nxt_M(x)=M^{-1}(s_M(x))$. A \emph{rotation exposed in $M$} is the cyclic sequence of pairs $(x_i,M(x_i))$ along a cycle $(x_0,\dots,x_{r-1})$ of $\nxt_M$, and eliminating it gives each $x_i$ the partner $M(x_{i+1})$. We use these facts.

\begin{enumerate}[label=(R\arabic*)]
\item Eliminating an exposed rotation from a stable matching gives a stable matching, and every stable matching is reached from $M_0$ by successive eliminations.
\item Let $\Pi$ be the set of rotations, ordered by precedence. Then $X\mapsto M_X$, where $M_X$ is $M_0$ with the rotations of $X$ eliminated in any order compatible with precedence, is a lattice isomorphism $\D(\Pi)\to\mathcal M_A$. The rotations exposed in $M_X$ are the minimal elements of $\Pi\setminus X$. Hence $\Pi\cong J(\mathcal M_A)$.
\item For an agent $a$, let $C_a\subseteq\Pi$ consist of the rotations whose pairs contain $a$. Then $C_a$ is a chain, and $M_X(a)$ depends only on $|X\cap C_a|$. Along $C_a$ a worker's partner strictly worsens and a firm's strictly improves.
\end{enumerate}

For (R3): eliminating a rotation changes exactly the partners of its members, a worker moving down its list and a firm up. If $\rho,\sigma\in C_a$ were incomparable, let $X$ be the union of their strict down-sets. Both would be minimal in $\Pi\setminus X$, hence exposed in $M_X$ by (R2), and both would contain $a$. But $\nxt_{M_X}$ is a function, so distinct exposed rotations share no worker, and hence no firm, since the firms of a rotation are the partners of its workers. Rotations outside $C_a$ do not move $a$, so $M_X(a)$ depends only on $X\cap C_a$, an initial segment of $C_a$ determined by its size.

\begin{enumerate}[label=(R\arabic*),resume]
\item \cite[Theorem~7]{GMRV} If $A,B$ differ in at most one worker's list (any number of firms may change), then $\mathcal M_A\cap\mathcal M_B$ is closed under the meet and join of $\mathcal M_A$.
\end{enumerate}

For completeness, here is a short proof of (R4). Let $M,M'$ be stable in both profiles. On each cycle of $M\cup M'$ all workers prefer the same one of $M,M'$ in a given profile: if $x$ prefers $y=M(x)$ to $M'(x)$, stability of $M'$ makes $y$ prefer $x_2=M'(y)$ to $x$, and stability of $M$ makes $x_2$ prefer $M(x_2)$ to $y$, and so on around the cycle. A cycle through $w$ also contains another worker, whose list is the same in $A$ and $B$. So the meets of $M,M'$ in $A$ and in $B$ coincide, and the meet lies in $\mathcal M_A\cap\mathcal M_B$. Joins are similar.

\section{Necessity}

Let $A,B$ differ at most in the lists of a worker $w$ and a firm $f$. If fewer lists differ, choose $w$ or $f$ arbitrarily. Use (R2) to identify $\mathcal M_A$ with $\D(\Pi)$, and let
\[
R^*=\{X\in\D(\Pi):\ M_X\in\mathcal M_B\}.
\]
Put $C_1=C_w=(\alpha_1<\dots<\alpha_p)$ and $C_2=C_f=(\beta_1<\dots<\beta_q)$, and write $\pi_1(X)=|X\cap C_1|$ and $\pi_2(X)=|X\cap C_2|$. Let $p_i$ be $w$'s partner when $\pi_1=i$, and $q_j$ be $f$'s partner when $\pi_2=j$. By (R3), $p_0\succ^A_w p_1\succ^A_w\cdots$ and $q_0\prec^A_f q_1\prec^A_f\cdots$. Write $\alpha_0=\top$ and $\beta_{q+1}=\bot$.

If $R^*=\varnothing$, take $A_1=\varnothing$, any $A_2$ and $\Phi=\varnothing$. Otherwise let $A_1=\pi_1(R^*)$ and $A_2=\pi_2(R^*)$, and let $\Phi$ be the set of all implications of either kind that every member of $R^*$ satisfies. Then $R^*\subseteq\Sol$. We prove $\Sol\subseteq R^*$.

\emph{Thresholds.} For a firm $g\ne f$, the partners of $g$ improve along $C_g$, so the down-sets $X$ with $w\succ^A_g M_X(g)$ are those with $|X\cap C_g|<\tau$ for some $\tau$. Accordingly there is $y_g\in C_g\cup\{\top,\bot\}$ with
\[
 w\succ^A_g M_X(g)\iff y_g\notin X\qquad(X\in\D(\Pi)).
\]
Symmetrically, for a worker $v\ne w$ there is $z_v\in C_v\cup\{\top,\bot\}$ with $f\succ^A_v M_X(v)\iff z_v\in X$. For $g=p_k\ne f$ with $k\ge1$ we have $y_{p_k}=\alpha_k$, because $p_k$ receives $w$ exactly when $\alpha_k$ is eliminated. Likewise $z_{q_l}=\beta_{l+1}$ for $l<q$ with $q_l\ne w$, because $q_l$ holds $f$ or a better firm until $\beta_{l+1}$ is eliminated, which moves it from $f$ to a worse firm.

\begin{lemma}\label{lem:order}
Let $i\in A_1$ and $i<k\le p$ with $p_k\ne f$. Then $p_i\succ^B_w p_k$. Let $j\in A_2$ and $0\le l<j$ with $q_l\ne w$. Then $q_j\succ^B_f q_l$.
\end{lemma}

\begin{proof}
Take $Y\in R^*$ with $\pi_1(Y)=i$. Then $\alpha_k\notin Y$, so $w\succ_{p_k}M_Y(p_k)$, and $p_k$'s list is unchanged. If $p_k\succ^B_w p_i$, the pair $(w,p_k)$ would block $M_Y$ in $B$. The second claim is symmetric, using $Y$ with $\pi_2(Y)=j$, so that $\beta_{l+1}\in Y$.
\end{proof}

\begin{lemma}\label{lem:couple}
Suppose $M_{X_0}(w)=f$ for some $X_0$. Then there are $j_0,j_1$ such that, for every $X\in\D(\Pi)$, $M_X(w)=f$ iff $\pi_1(X)=j_0$ iff $\pi_2(X)=j_1$, and $\pi_1(X)\ge j_0$ iff $\pi_2(X)\ge j_1$.
\end{lemma}

\begin{proof}
Let $f=p_{j_0}$ and $w=q_{j_1}$. If $j_0=0$ then $w,f$ are partners in $M_0$, so $j_1=0$ and the last claim is trivial. If $j_0\ge1$, the rotation $\alpha_{j_0}$ moves $w$ onto $f$, so it changes $f$'s partner to $w$ and equals $\beta_{j_1}$. Thus $\pi_1(X)\ge j_0\iff\alpha_{j_0}\in X\iff\pi_2(X)\ge j_1$.
\end{proof}

\begin{proof}[Proof of $\Sol\subseteq R^*$]
Let $X\in\Sol$, $i=\pi_1(X)$ and $j=\pi_2(X)$, and suppose a pair blocks $M_X$ in $B$. Pairs avoiding $w$ and $f$ have their $A$-preferences, and $M_X$ is $A$-stable, so the pair contains $w$ or $f$. In each case below we exhibit an implication of $\Phi$'s kind that $X$ violates, and derive a contradiction from its failure on $R^*$; hence it lies in $\Phi$, contradicting $X\in\Sol$.

\emph{(a) The pair is $(w,g)$ with $g\ne f$.} Then $g\succ^B_w p_i$ and $y_g\notin X$, so $y_g\ne\top$. The implication $\alpha_i\to y_g$ is of $w$-kind and $X$ violates it. Suppose $Y\in R^*$ violates it too, and take $Y'\in R^*$ with $\pi_1(Y')=i$. By (R4), $Z=Y\cap Y'\in R^*$. Since $\alpha_i\in Y$ we get $\pi_1(Z)=i$, and $y_g\notin Z$. So $(w,g)$ blocks $M_Z$ in $B$, a contradiction.

\emph{(b) The pair is $(v,f)$ with $v\ne w$.} Then $v\succ^B_f q_j$ and $z_v\in X$, so $z_v\ne\bot$. The implication $z_v\to\beta_{j+1}$ is of $f$-kind and $X$ violates it. If $Y\in R^*$ violates it, take $Y'\in R^*$ with $\pi_2(Y')=j$ and $Z=Y\cup Y'$. Then $\pi_2(Z)=j$ and $z_v\in Z$, so $(v,f)$ blocks $M_Z$ in $B$.

\emph{(c) The pair is $(w,f)$.} Then $M_X(w)\ne f$, $f\succ^B_w p_i$ and $w\succ^B_f q_j$. The implication $\alpha_i\to\beta_{j+1}$ is of both kinds and $X$ violates it. Let $Y\in R^*$ violate it, with $k=\pi_1(Y)\ge i$ and $l=\pi_2(Y)\le j$. If $M_Y(w)\ne f$, then $p_k\ne f$ and $q_l\ne w$, and Lemma~\ref{lem:order} gives $p_i\succeq^B_w p_k$ and $q_j\succeq^B_f q_l$ (the first trivially if $k=i$, the second if $l=j$). Hence $f\succ^B_w p_k$ and $w\succ^B_f q_l$, and $(w,f)$ blocks $M_Y$. If $M_Y(w)=f$, Lemma~\ref{lem:couple} gives $k=j_0$ and $l=j_1$. Then $\pi_2(X)=j\ge j_1$ forces $i\ge j_0=k\ge i$, so $\pi_1(X)=j_0$ and $M_X(w)=f$, a contradiction.
\end{proof}

So $R^*=\Sol$. A lattice isomorphism $\mathcal M_A\to L$ carrying $\mathcal M_A\cap\mathcal M_B$ onto $R$ restricts to a poset isomorphism $\Pi\cong J(\mathcal M_A)\to J(L)=P$. It carries our system on $\Pi$ to a one-change system on $P$ whose solution set corresponds to $R$. This proves necessity.

\section{Sufficiency}

Fix a finite poset $P$ and a one-change system on it with chains $C_1=(\alpha_1<\dots<\alpha_p)$ and $C_2=(\beta_1<\dots<\beta_q)$. We build $A$ and $B$.

\emph{Agents.} The workers are $m_\rho,m'_\rho$ for $\rho\in P$, together with $w,c_1,\dots,c_p,d_0,d_1,\dots,d_q,u,v$. The firms are $h_\rho,h'_\rho$ for $\rho\in P$, together with $e_0,\dots,e_p,f,k_1,\dots,k_q,g,\hat f$.

\emph{Cycles.} For $\rho\in P$ let $\Gamma_\rho$ be the cyclic sequence of workers formed by concatenating $(w,c_t)$ if $\rho=\alpha_t$, then $(d_s,d_{s-1})$ if $\rho=\beta_s$, then $(m_\rho,m'_\rho)$. The initial matching is
\[
M_\varnothing=\{m_\rho h_\rho,\ m'_\rho h'_\rho\ (\rho\in P),\ we_0,\ c_te_t,\ d_0f,\ d_sk_s\ (s\ge1),\ ug,\ v\hat f\}.
\]
For $X\in\D(P)$, the matching $M_X$ is obtained from $M_\varnothing$ by \emph{eliminating} the members of $X$ in an order compatible with $P$, where eliminating $\rho$ gives each worker of $\Gamma_\rho$ the current partner of the next worker of $\Gamma_\rho$. Two cycles share a worker only if both rotations lie in $C_1$ or both lie in $C_2$, hence only if they are comparable. Eliminations along disjoint cycles commute, and any two orders compatible with $P$ differ by swaps of adjacent incomparable elements, so $M_X$ does not depend on the order chosen.

Each $m_\rho,m'_\rho$ moves only at $\rho$, each $c_t$ only at $\alpha_t$, each $d_s$ only at $\beta_s$ and $\beta_{s+1}$, and $w$ at $\alpha_1,\dots,\alpha_p$ in turn. So each agent's partners are as follows. Put $\kappa_t=k_s$ if $\alpha_t=\beta_s$, and $\kappa_t=h_{\alpha_t}$ otherwise. Put $\varepsilon_\rho=e_{t-1}$ if $\rho=\alpha_t$; otherwise $\varepsilon_\rho=k_s$ if $\rho=\beta_s$; otherwise $\varepsilon_\rho=h_\rho$. Put $\lambda_\rho=d_{s-1}$ if $\rho=\beta_s$; otherwise $\lambda_\rho=c_t$ if $\rho=\alpha_t$; otherwise $\lambda_\rho=m'_\rho$.
\begin{itemize}
\item Workers, in order of occurrence: $w$: $e_0,e_1,\dots,e_p$. $c_t$: $e_t,\kappa_t$. $d_s$: $k_s$ (if $s\ge1$), $f$, $h_{\beta_{s+1}}$ (if $s<q$). $m_\rho$: $h_\rho,h'_\rho$. $m'_\rho$: $h'_\rho,\varepsilon_\rho$. $u$: $g$. $v$: $\hat f$.
\item Firms, in order of occurrence: $e_t$: $c_t$ (if $t\ge1$), $w$, $m'_{\alpha_{t+1}}$ (if $t<p$). $f$: $d_0,\dots,d_q$. $k_s$: $d_s$, then $c_t$ if $\beta_s=\alpha_t$ and $m'_{\beta_s}$ otherwise. $h_\rho$: $m_\rho,\lambda_\rho$. $h'_\rho$: $m'_\rho,m_\rho$. $g$: $u$. $\hat f$: $v$.
\end{itemize}

\emph{Profile $A$.} A worker lists its partners in order of occurrence and a firm in reverse order, so each agent prefers the later partners if it is a firm and the earlier ones if it is a worker. Then make these insertions.
\begin{enumerate}[label=(I\arabic*)]
\item For each covering pair $\sigma\lessdot\tau$ of $P$: $h_\sigma$ into $m_\tau$'s list directly after $h_\tau$, and $m_\tau$ into $h_\sigma$'s list directly before $m_\sigma$.
\item For each $y\in P$ that is the head of a $w$-kind implication: $w$ into $h_y$'s list directly before $m_y$.
\item For each $z\in P$ that is the tail of an $f$-kind implication: $f$ into $m_z$'s list directly after $h_z$.
\item $g$'s list begins $w,u$, and $v$'s list begins $f,\hat f$.
\end{enumerate}
Insertions at the same place are made in any order. All remaining agents follow, in any order, at the end of every list. Call these the \emph{tails} of the lists.

\begin{lemma}\label{lem:A}
The stable matchings of $A$ are exactly the $M_X$ with $X\in\D(P)$. They are distinct, and $X\mapsto M_X$ is a lattice isomorphism $\D(P)\to\mathcal M_A$. Moreover $M_X(w)=e_{|X\cap C_1|}$ and $M_X(f)=d_{|X\cap C_2|}$; $h_y$ holds $m_y$ iff $y\notin X$; and $m_z$ holds $h_z$ iff $z\notin X$.
\end{lemma}

\begin{proof}
In worker-proposing deferred acceptance every worker first proposes to its first choice. These are distinct except that $d_0$ and $v$ both propose to $f$, which keeps $d_0$ because $v$ lies in its tail. Then $v$ is accepted by $\hat f$. So $M_\varnothing$ is the worker-optimal stable matching.

Fix $X\in\D(P)$. A worker is \emph{finished} if it holds its last partner. For an unfinished worker $x$, let $\cur(x)$ be the least member of $C_x\setminus X$, where now $C_x=\{\rho:x\in\Gamma_\rho\}$; this is the rotation that next moves $x$. We compute $\nxt=\nxt_{M_X}$. A firm prefers $x$ to its partner only if $x$ is above that partner on its list. Before its tail, every firm lists only its partners, the insertions (I1)--(I4), and nothing else, and every firm lists its last partner first, except $g$. Using this, the successors are as follows.
\begin{enumerate}[label=(\roman*)]
\item $x=m_\tau$, $\tau\notin X$. It holds $h_\tau$. Next come, in some order, the firms $h_\sigma$ for $\sigma\lessdot\tau$ and possibly $f$, then $h'_\tau$. The firm $f$ keeps $m_\tau$ in its tail. The firm $h_\sigma$ prefers $m_\tau$ iff it holds $m_\sigma$, i.e.\ iff $\sigma\notin X$. If some such $\sigma\notin X$, then $\nxt(x)=m_\sigma$ for the first such $\sigma$ on the list, and $\cur(m_\sigma)=\sigma<\tau$. Otherwise $\nxt(x)=m'_\tau$, and every lower cover of $\tau$, hence everything below $\tau$, lies in $X$.
\item $x=m'_\tau$, $\tau\notin X$. It holds $h'_\tau$ and its next entry is $\varepsilon_\tau$, which lists $m'_\tau$ first. If $\tau=\alpha_t$, the holder of $e_{t-1}$ is $w$ when $\alpha_{t-1}\in X$ (or $t=1$), with $\cur(w)=\tau$; otherwise it is $c_{t-1}$, with $\cur(c_{t-1})=\alpha_{t-1}<\tau$. If $\tau=\beta_s\notin C_1$, the holder of $k_s$ is $d_s$, with $\cur(d_s)=\tau$. Otherwise the holder of $h_\tau$ is $m_\tau$.
\item $x=w$ holding $e_i$ with $i<p$. Its next entry $e_{i+1}$ holds $c_{i+1}$ and prefers $w$, so $\nxt(w)=c_{i+1}$ and $\cur(c_{i+1})=\alpha_{i+1}=\cur(w)$.
\item $x=c_t$, $\alpha_t\notin X$. Its next entry $\kappa_t$ lists $c_t$ first and holds $m_{\alpha_t}$ or $d_s$, both with $\cur=\alpha_t$.
\item $x=d_s$ with $s\ge1$ and $\beta_s\notin X$. Its next entry is $f$, which holds $d_j$ with $j<s$ and prefers $d_s$. If $j=s-1$ then $\cur(d_{s-1})=\beta_s$; otherwise $\cur(d_j)=\beta_{j+1}<\beta_s$.
\item $x=d_s$ holding $f$ with $s<q$. Its next entry $h_{\beta_{s+1}}$ lists $d_s$ first and holds $m_{\beta_{s+1}}$, with $\cur=\beta_{s+1}=\cur(d_s)$.
\item Finished workers. A finished worker $x\ne w$ has only its tail after its partner. A firm there is not a partner of $x$, and the only insertions of such workers into firm lists, (I1) and $u$ in (I4), go into firms that precede their tails. So the firm lists $x$ in its own tail, below its partner, and $x$ has no successor. The worker $w$ holding $e_p$ may have successor $g$, held by $u$, or some $h_y$ with $y\notin X$, held by $m_y$. No worker has $w$ as its $\nxt$, since $e_p$ lists $w$ first.
\end{enumerate}
In each case where $\cur$ is unchanged, the new worker follows $x$ in the cyclic order of $\Gamma_{\cur(x)}$. Thus every edge of $\nxt$ either lowers $\cur$ strictly, or joins consecutive workers of $\Gamma_{\cur(x)}$, or leaves $w$ for $u$ or for an unfinished worker. No cycle contains a finished worker, since by (vii) only $w$ among them has a successor and no worker has $w$ as its $\nxt$. So along a cycle $\cur$ is constant, say $\tau$, and the cycle is $\Gamma_\tau$. By (i) this forces $\tau$ to be minimal in $P\setminus X$. Conversely, if $\tau$ is minimal, then (i)--(vi) supply every edge of $\Gamma_\tau$: when $\tau=\alpha_t$ we have $\alpha_{t-1}\in X$, and when $\tau=\beta_s$ the firm $f$ holds $d_{s-1}$.

So the rotations exposed in $M_X$ are the $\Gamma_\tau$ with $\tau$ minimal in $P\setminus X$, and eliminating $\Gamma_\tau$ yields $M_{X\cup\{\tau\}}$. By (R1) and induction on $|X|$, every $M_X$ is stable and these are all the stable matchings. The set $X$ is recovered from $M_X$ as $\{\tau: m_\tau\text{ holds }h'_\tau\}$. Workers only lose by eliminations, so $X\subseteq Y$ implies $M_X\le M_Y$. Conversely, if $M_X\le M_Y$ and $\tau\in X$, then $m_\tau$ holds its last partner $h'_\tau$ in $M_X$, hence in $M_Y$, so $\tau\in Y$. The remaining claims are read off the partner lists.
\end{proof}

\emph{Profile $B$.} Only $w$ and $f$ change. Each agent below is listed only at its first occurrence. With $\alpha_0=\top$, the list of $w$ is: for $i=0,1,\dots,p$ in turn, first the firms $h_y$ for each $(\alpha_i\to y)\in\Phi$ with $y\in P$, and $g$ if $(\alpha_i\to\bot)\in\Phi$, then $e_i$ if $i\in A_1$. Next comes $g$ if $A_1\ne\{0,\dots,p\}$, then the $e_i$ with $i\notin A_1$, then everyone else. The list of $f$ is: first the workers $m_z$ for each $(z\to\bot)\in\Phi$ with $z\in P$, and $v$ if $(\top\to\bot)\in\Phi$. Then, for $l=q,q-1,\dots,0$ in turn: $d_l$ if $l\in A_2$, followed, when $l\ge1$, by the workers $m_z$ for each $(z\to\beta_l)\in\Phi$ with $z\in P$, and $v$ if $(\top\to\beta_l)\in\Phi$. Next comes $v$ if $A_2\ne\{0,\dots,q\}$, then the $d_l$ with $l\notin A_2$, then everyone else.

Two implications with the same head $y$ and tails $\alpha_k,\alpha_{k'}$, $k<k'$, are together equivalent to the one with tail $\alpha_k$, since $\alpha_{k'}\in X$ implies $\alpha_k\in X$. Symmetrically on the $f$ side, the head $\beta_l$ with the largest $l$ is the binding one. So listing each agent only at its first occurrence realizes all of $\Phi$.

\begin{lemma}\label{lem:B}
For $X\in\D(P)$, $M_X$ is stable in $B$ iff $X\in\Sol$.
\end{lemma}

\begin{proof}
Let $i=|X\cap C_1|$ and $j=|X\cap C_2|$, so $M_X(w)=e_i$ and $M_X(f)=d_j$. Pairs avoiding $w,f$ keep their $A$-lists and do not block. The pair $(w,f)$ does not block, since $f$ lies after $e_i$ on $w$'s new list.

Pairs $(w,y)$ with $y\ne f$ above $e_i$ on $w$'s new list. If $i\notin A_1$, then $g$ is such a firm and always prefers $w$ to $u$, so $(w,g)$ blocks; and $X\notin\Sol$. If $i\in A_1$, the firms above $e_i$ are of three sorts. The first is $e_k$ with $k<i$; it holds $m'_{\alpha_{k+1}}$, its first choice, so it does not block. The second is $h_y$ listed for $(\alpha_k\to y)$ with $k\le i$; by Lemma~\ref{lem:A} and (I2) it prefers $w$ iff $y\notin X$, so it blocks iff that implication fails, since $\alpha_k\in X$. The third is $g$ listed for $(\alpha_k\to\bot)$ with $k\le i$; it blocks, matching the failure of that implication.

Pairs $(x,f)$ with $x\ne w$ above $d_j$ on $f$'s new list. If $j\notin A_2$, then $v$ is such a worker, always prefers $f$ to $\hat f$, and blocks. If $j\in A_2$, the workers above $d_j$ are of three sorts. The first is $d_l$ with $l>j$; it holds $k_l$, its first choice, so it does not block. The second is $m_z$ listed for $(z\to\beta_l)$ with $l>j$, or for $(z\to\bot)$; by Lemma~\ref{lem:A} and (I3), $m_z$ prefers $f$ to its partner iff $z\in X$, so it blocks iff the implication fails, since $\beta_l\notin X$. The third is $v$ listed for $(\top\to\beta_l)$ with $l>j$, or for $(\top\to\bot)$, which blocks, matching the failure of that implication.

Hence $M_X$ is $B$-stable iff $i\in A_1$, $j\in A_2$ and $X$ satisfies $\Phi$.
\end{proof}

Lemmas~\ref{lem:A} and~\ref{lem:B} prove sufficiency: $A$ and $B$ differ only in the lists of $w$ and $f$, and the isomorphism $X\mapsto M_X$ carries $\Sol$ onto $\mathcal M_A\cap\mathcal M_B$. With necessity, this proves Theorem~\ref{thm:main}.

\begin{remark}
As a check on the construction, a program built $A$ and $B$ for 3{,}800 random systems on posets with up to six elements. In every case the stable matchings of $A$ were exactly the $M_X$, and $\mathcal M_A\cap\mathcal M_B$ was exactly $\Sol$. Independently, exhaustive enumeration of single worker and firm changes on random instances realized all $317$ isomorphism types of pairs $(L,R)$ with $|L|\le6$ and $R$ empty or a sublattice. These computations are not used in the proofs.
\end{remark}

\section*{Completion Estimate}
100\%

The claim covers the whole catalogue target: necessary and sufficient conditions on $(L,R)$, including the empty case, and the closure behaviour of $L\setminus R$. It has not been independently refereed.

\begin{thebibliography}{9}
\bibitem{GS} D. Gale and L. S. Shapley, College admissions and the stability of marriage, \emph{Amer. Math. Monthly} 69 (1962), 9--15.
\bibitem{GI} D. Gusfield and R. W. Irving, \emph{The Stable Marriage Problem: Structure and Algorithms}, MIT Press, 1989.
\bibitem{GMRV} R. R. Gangam, T. Mai, N. Raju and V. V. Vazirani, Stable matching: dealing with changes in preferences, arXiv:2304.02590v3 (2025), Theorem~7 and Section~4.
\end{thebibliography}

\end{document}
